\section{Pointwise convergence for the proper subsets of four times}
\label{sec:lower-length}

We prove Lemma~\ref{lem:lower-length}.  The systems to which that lemma
is applied in Section~\ref{sec:selected} are extensions of a suspension
and need not be mixing or ergodic.  Accordingly, the input needed here
is a pointwise convergence theorem on arbitrary invertible probability
preserving systems.  Its analytic basis is the oscillation estimate in
\cite[Proposition~2.2]{TripleCompanion}, denoted by P in this appendix.  We also
write H for \cite{HilbertCompanion}.

P states this oscillation estimate for the times $n,2n,3n$ before using
mixing to identify its limit.  We need the three additional triples
contained in $\{1,2,3,4\}$.  We obtain them by modifying the local
four-linear construction in H and its extension to inputs varying with
depth in P.  This is a use of those analytic proofs, including their
intermediate estimates; the boundedness theorem for the trilinear
Hilbert transform alone would not supply the required oscillation
estimate.  We specify the imported statements first, then verify the
coefficient changes, and finally give the passage to arbitrary systems
and all integer averaging lengths.

\subsection{The local estimates used from H and P}

For this appendix, a \emph{role set} is a set $S$ of four distinct
integer coefficients.  Initially $S=\{0\}\cup J$, where
$J\subset\{1,2,3,4\}$ has three elements.  If $K$ is a dyadic interval
of length $r$, put
\begin{equation}
 \mathcal H_{K,S}(z)
 =r^{-2}\int_{\R^2}w_K(x,t)\prod_{j\in S}z_j(x+jt)\,\dd x\,\dd t.
 \label{eq:ll-local-form}
\end{equation}
Fix $c_0>0$ and $C_0\ge1$.  The admissible kernels are the smooth,
compactly supported functions $w_K$ for which every $x+jt$ on the
support lies in $K$ at distance at least $c_0r$ from its boundary,
\begin{equation}
 \int_{\R}w_K(v-jt,t)\,\dd t=0
 \quad(v\in\R, j\in S),
 \qquad
 \|\partial_x^a\partial_t^b w_K\|_\infty
 \le r^{-a-b}\bigl(C_0(a+b+2)\bigr)^{C_0(a+b+2)}.
 \label{eq:ll-kernels}
\end{equation}
Kernels may be chosen independently at the different intervals.

H's \emph{Uniform local estimate}
\cite[Theorem~2.2]{HilbertCompanion}
asserts, for $S=\{0,1,2,3\}$ and a suitable $q\in(2,3)$, that
\begin{equation}
 \sum_{K\in\mathcal K}|K|\,|\mathcal H_{K,S}(z)|
 \le C\int_{\R}\prod_{j\in S}M_qz_j(y)\,\dd y,
 \qquad M_qz=(M(|z|^q))^{1/q},
 \label{eq:ll-fixed-local}
\end{equation}
for every finite family of dyadic intervals and compactly supported
scalar $L^q$ inputs.  Here $M$ is the Hardy--Littlewood maximal
operator.  Its proof, in H, Sections~3--10, supplies the quantitative
chart, projection, separation, counting, and side estimates used in P.

The extension needed for oscillation is P's \emph{Local estimate with
variation in depth}, Theorem~3.2.  To state it, let $K_*$ be a root
interval and let $s(K)=\log_2(|K_*|/|K|)$ be the depth of a descendant.
For a finite integer block $B=\{a,\ldots,b\}$ and a scalar sequence
$(z_s)$, write
\[
 \mathcal V_B(z)(y)
 =\max_{s\in B}|z_s(y)|+
   \sum_{s=a}^{b-1}|z_{s+1}(y)-z_s(y)|.
\]
For the original role set $\{0,1,2,3\}$, P proves the assertion below.
We establish it for the four required role sets.

\begin{proposition}[The local analytic input with changed coefficients]
\label{prop:ll-depth}
Let $S=\{0\}\cup J$, where $J$ is a three-element subset of
$\{1,2,3,4\}$.  Let $\mathcal T$ be a finite rooted dyadic tree with
root $K_*$ and largest depth $d$, and let $\mathcal K\subset\mathcal T$
have kernels satisfying \eqref{eq:ll-kernels}.  For each $j\in S$, let
$(z_{j,s})_{0\le s\le d}$ be measurable scalar functions on $K_*$.
Extend these functions by zero outside $K_*$.
Suppose that $\{0,\ldots,d\}$ has a deterministic partition into
$W_j\ge1$ nonempty consecutive blocks, on each of which
\[
 \mathcal V_B(z_j)(y)\le W_j^{-1/2}
 \quad\text{for almost every }y\in K_*.
\]
The partitions may differ between the four roles.  Then
\begin{equation}
 \sum_{K\in\mathcal K}|K|
   \bigl|\mathcal H_{K,S}((z_{j,s(K)})_{j\in S})\bigr|
 \le C(c_0,C_0)|K_*|.
 \label{eq:ll-depth}
\end{equation}
The constant is independent of the tree, kernels, depth range, and
block partitions.  It can be chosen uniformly over the four sets $S$.
\end{proposition}

The proof of Proposition~\ref{prop:ll-depth} is the proof-by-modification
below.  The analytic inputs are H's local estimate and its proofs of the
intermediate estimates, and P's depth-local construction in Sections~4--7.
Proposition~\ref{prop:ll-depth} is the sole analytic input to the
local-to-orbit argument in Subsection~\ref{subsec:ll-transfer};
Subsection~\ref{subsec:ll-convergence} then proves
Lemma~\ref{lem:lower-length}.  The intervening subsections establish
the proposition by checking, in order, the coordinate and quadratic
identities, the integer-lift and counting arithmetic, and their
use in P's depth scheduling.
We retain their definitions of quadratic chart atoms, inherited
projection families, curvature tags, and finite progression tests.  The
following computations specify their dependence on the four role
coefficients.  In particular, the proposition is not an assertion of
uniform continuity of a Hilbert-transform bound in its slopes.

\subsection{Coordinate changes and the quadratic cancellation vector}
\label{subsec:ll-algebra}

Every set $S$ above contains an adjacent pair.  Translate all role
labels by the same integer so that one such pair is $0,1$.  This is the
change of variables $x'=x+mt$ in \eqref{eq:ll-local-form}, with
Jacobian one.  The four spatial coordinates themselves are unchanged.
Consequently the support and cancellation conditions remain the same,
and the derivative bounds retain their form after increasing the fixed
constant $C_0$.  During the local argument we use these translated
labels.  Their pairwise differences have absolute value at most four,
and their diameter is either three or four.  Thus
\begin{equation}
 |t|\le |K|/3
 \quad\text{whenever every }x+jt\text{ lies in }K.
 \label{eq:ll-support-t}
\end{equation}

The vector $(-1,3,-3,1)$ in H is replaced by a primitive integer vector
$(c_j)_{j\in S}$ proportional to
\[
 \left(\prod_{\ell\in S\setminus\{j\}}(j-\ell)\right)^{-1}.
\]
The overall sign is immaterial.  Before translation, convenient choices
are
\begin{equation}
\begin{array}{c|rrrr}
 S&\multicolumn{4}{c}{(c_j)_{j\in S}\text{ in increasing role order}}\\ \hline
 \{0,1,2,3\}&-1&3&-3&1\\
 \{0,1,2,4\}&-3&8&-6&1\\
 \{0,1,3,4\}&-1&2&-2&1\\
 \{0,2,3,4\}&-1&6&-8&3
\end{array}
\label{eq:ll-coefficient-table}
\end{equation}
Lagrange interpolation of the polynomials $1,u,u^2$ gives
\begin{equation}
 \sum_{j\in S}c_j=\sum_{j\in S}jc_j=
 \sum_{j\in S}j^2c_j=0.
 \label{eq:ll-moments}
\end{equation}
These identities are preserved by translation of the role labels.
Every $c_j$ is nonzero.  All coefficients and denominators arising
from this finite list of coordinate systems are bounded absolute
constants.

We distinguish numerical role labels from enumerations of four
objects.  A coefficient such as $j$, $j-i$, or $c_j$ is changed according
to $S$.  An ordering of four input sizes, four leaves, or four positions
in a multilinear expansion still enumerates four objects.  This
convention leaves those parts of H's proofs unchanged.

\subsection{Chart separation, detection, and the scale estimates}

The local absolute estimate, H, Lemma~2.1, uses the map
$(x,t)\mapsto(x+it,x+jt)$, whose determinant is $j-i$.  Its endpoint
bounds and interpolation therefore hold for the new roles.  In
particular the initial bound for the finite-depth constant in
\eqref{eq:ll-fixed-local}, and H's stopping reduction to normalized
local sizes, remain valid.  The coefficient and projection bounds,
compression and height-bin estimates, H, Lemmas~3.2--3.5, are
statements about finite families of vectors or real frequencies and do
not specify a progression coefficient.

We next check the part of H, Section~4, that converts two chart atoms
in distinct roles into separated factors.  A chart phase has the form
\[
 \lambda y^2+(\theta y)^TCd+\beta y,
\]
where $\theta$ is a vector of real frequencies, $C$ is the rational
chart matrix, and $d$ is a bounded lift of $\theta y$ modulo the
integer lattice.  The source definition also includes a smooth
amplitude on the lift.  Its one-variable Weyl estimates, Fourier
expansion, and uniform tail estimates are unchanged.

For H, Lemma~4.5 (\emph{Progression separation}), fix source roles $i,k$
and spectator roles $a,b$.  Set $\alpha_i=1$, $\alpha_k=0$, and solve
\[
 \sum_{j\in S}\alpha_j=\sum_{j\in S}j\alpha_j=0
\]
for $\alpha_a,\alpha_b$.  This is possible because $a\ne b$.
Writing $y_j=x+jt$, insert the same lift partitions as in H at
$\theta y_a$, $\theta y_b$, and $\theta t$.  If their lifts are
$d_a,d_b,d$, then
\[
 k_j=d_j-d_a-(j-a)d
\]
is an integer vector.  For instance, if $d_j=\theta y_j-m_j$ and
$d=\theta t-m$, the identity $y_j=y_a+(j-a)t$ gives
$k_j=-m_j+m_a+(j-a)m$.  Thus the equality cutoffs in H still test
integer carries exactly.  Their support bounds change only by a
fixed factor.  With $A_2=\sum_jj^2\alpha_j$, direct expansion gives
\begin{equation}
 \begin{split}
 \sum_{j\in S}\alpha_j
       [\lambda y_j^2+(\theta y_j)^TCd_j]
 &=A_2[\lambda t^2+(\theta t)^TCd]\\
 &\quad+\sum_{j\in S}\alpha_j(\theta y_j)^TCk_j.
 \end{split}
 \label{eq:ll-chart-identity}
\end{equation}
The last sum is linear in $x,t$ and hence separates in $y_a,y_b$.
The smooth amplitude on the bounded lift box has the same Gevrey
bounds up to fixed constants.  Fourier expansion therefore gives
H's absolutely convergent separation into two spectator atoms and one
atom in $t$, with the same coefficient-tail exponents and
quasipolynomial complexity bounds.  H, Lemma~4.6 (\emph{Single-scale separation}), then uses $t=(y_b-y_a)/(b-a)$; every added rational height is
bounded by an absolute constant.

For the generalized von Neumann bound, H, Lemma~4.11, let $i$ be
the role whose input is controlled by its $U^3$ norm.
Apply Cauchy--Schwarz in the directions $(-j,1)$ for the other
three roles, beginning with the role containing the $L^2$ input.
Each direction holds $x+jt$ fixed and changes $x+it$ by
$(i-j)h$.  Multiplication by the nonzero integer $i-j$ preserves Haar
measure on a circle.  The three independent circle parameters
therefore produce the same cube integral as in the source proof.
Character weights split using the marked pair $x=y_0$,
$t=y_1-y_0$.  Enlarging the fixed periodic boxes accommodates all
translated role sets.  This verifies the coefficient-dependent part
of the generalized von Neumann bound.  In H, Lemma~4.12
(\emph{Local detection}), the potentially unbounded $L^2$ residual
occupies a different role.  After clipping the opposing inputs, the
source proof decomposes two of them using its continuous inverse
theorem and bounds the uniform errors by Lemma~4.11.  The separated
chart moments then detect the residual.  Those separations were
verified above; the clipping, inverse theorem, and mixed-norm
exponents are unchanged.

H, Lemma~5.1 (\emph{A scale sum for one separated term}),
uses the spectator coordinates $(y_a,y_b)$.  Holding one fixed, the
integral of the kernel over the other is zero by
\eqref{eq:ll-kernels}, with a constant Jacobian factor $|b-a|$.
The double primitive of that kernel consequently has the same
compact support and smooth bounds.  Its tensor expansion into
mean-zero bumps, the Bessel estimates for those bumps, and the
frequency-height bounds apply with the new pair.  The only additional
demodulations have frequencies $\pm\gamma/(b-a)$; their rational
multipliers have bounded height, although $\gamma$ can be arbitrary
real.  This also verifies the coefficient changes
in the finite predictor list and the delayed predictor replacement.
The scale bands, accuracy gaps, and projection-energy bounds in the
remaining base construction depend on these estimates and the
four-slot multilinearity, rather than the values of the labels.
There is also a multivariable use of the chart estimates in H,
Lemma~5.7 (\emph{Few active scales for a fixed triple}).  Squaring
the transpose at the remaining role $\ell$ gives six atoms with
arguments $h+(j-\ell)t$ and $h+(j-\ell)t'$.  These are still fixed
integer linear forms in three variables, so the three-variable Weyl
alternative and chart expansion of H, Lemmas~4.3--4.4, apply with
the same complexity bounds.  Its final kernel cancellation is
the marginal at anchor $\ell$ in \eqref{eq:ll-kernels}.

\subsection{Curvature tags, short blocks, and integer heights}

H's localization in Section~6 groups chart atoms by their normalized
quadratic phase and horizontal chart data, called their curvature
tags.  Normalize the parameters $(\lambda,C)$ in role $j$ by
dividing both by $c_j$, and retain the horizontal vector $\theta$.
The coefficient calculation in H, Lemma~6.1 (\emph{Separation of different tags}), can
then be checked directly.  For source roles $i,j$ and spectator roles
$a,b$, put
\[
 \xi_\ell=\frac{b-\ell}{b-a},\qquad
 \zeta_\ell=\frac{\ell-a}{b-a};
 \qquad y_\ell=\xi_\ell y_a+\zeta_\ell y_b.
\]
The barycentric formula gives
\begin{equation}
 c_i\xi_i\zeta_i=-c_j\xi_j\zeta_j\ne0.
 \label{eq:ll-tag-identity}
\end{equation}
Indeed the factors $(i-a)(i-b)$ cancel from the product defining
$c_i$, leaving a nonzero multiple of $(i-j)^{-1}$; interchanging
$i,j$ changes its sign.  The return calculation in that lemma
therefore has the same nonzero prefactor and the same separation
conclusion.  Its period is enlarged to clear these fixed denominators.

The residue-template calculation in H, Lemma~6.6
(\emph{A uniform same-tag test}), uses the
marked adjacent pair.  Its lifts satisfy
\begin{equation}
 d_\ell=(1-\ell)d_0+\ell d_1+k_\ell,
 \qquad k_\ell\in\Z^d.
 \label{eq:ll-template-lift}
\end{equation}
Substitution of \eqref{eq:ll-template-lift} and
\eqref{eq:ll-moments} cancels the common quadratic terms and leaves
only linear phases from the carries.  Their Fourier expansion has
the same residue-template bound after a fixed enlargement of the
lift box.  It introduces no new curvature tag.  Thus the source
localization proof receives the same tag-separation and transfer
estimates.  Its graph distances, color histories, masks, and
projection bounds are then used with those estimates and their
original parameter order.

H, Lemma~7.2 (\emph{Short-block comparison}), uses two sampling
windows to compare a finite lattice form with a form on a connected
torus.  Place those windows on roles $0,1$.  The lattice and lifted
torus coordinates in the other roles are respectively
\[
 n_j=(1-j)n_0+jn_1,
 \qquad w_j=(1-j)u+jv.
\]
Both formulas have integer coefficients.  A fixed torus of length
$64$ and lift boxes in $[-12,12]$ accommodate the present role sets;
the source support parameter is decreased by a fixed factor if
necessary.  Maps from a pair of variables to two distinct role
coordinates are injective on their lattice images and surjective on
connected tori.  These are the pair-coordinate properties needed in
the comparison.

The quadratic randomization in role $j$ becomes
$e(\gamma m^2/c_j)$, where $e(v)=\exp(2\pi i v)$ and $\gamma$ is
uniform modulo a common multiple of all $|c_j|$; the value $24$ works
for \eqref{eq:ll-coefficient-table}.  On the Fourier support of the
four-linear form,
\[
 \sum_jm_j=\sum_jjm_j=0.
\]
The two vectors $c=(c_j)$ and $(jc_j)$ are independent and span this
two-dimensional nullspace.  Hence $m_j=c_j(v+jv')$ for some real
$v,v'$, and
\[
 \sum_j\frac{m_j^2}{c_j}
 =\sum_jc_j(v+jv')^2=0
\]
by \eqref{eq:ll-moments}.  The randomization preserves the form.
For each single role, averaging its fourth moment in the spatial
variable and in $\gamma$ imposes equality of a pair's sum and sum of
squares; the same two pairings as in H remain.  Thus its fourth-moment
and clipping estimates, including their powers of $q-2$, are
unchanged.

The return from lattice points to small intervals also holds with
these coordinates.  If four real role coordinates are within $e$ of
integers, the marked adjacent pair reconstructs an integer progression.
The discrepancy at any other role is an integer of magnitude less
than $8e$.  Choosing $e<1/100$ forces it to vanish.  The admissible
real offsets have positive area, computed explicitly below in
\eqref{eq:ll-polygon}.  All support and overlap losses in the
short-block comparison are therefore fixed constants.

The remaining arithmetic change is the height base.  In H,
Sections~8--9, replace $L_h=6^h$ by $L_h=12^h$, and replace the
factor $36$ in the finite progression tests by $144$.  Anchor changes
still have integer coefficients.  The common progression multiplier
in H, Section~8.3, uses exactly \eqref{eq:ll-template-lift} and
\eqref{eq:ll-moments}, so its smoothness and Fourier truncation are
unchanged in type.  For the divisibility step, a coordinate in the
free lattice quotient in H, Lemma~8.3 (\emph{Height stabilization}),
is an integral rational multiple of $12^{R'}$.  Here the quotient
is the integer coordinate lattice modulo its intersection with the
source's small rational subspace, $R'$ is the height gap, and $C_*$
bounds the denominators before multiplication by $12^{R'}$.
The intersection is saturated, so this quotient is free abelian.
Writing $\Delta$ for the required drop in height, its $2$- and $3$-adic
valuations are at least
\[
 2R'-\log_2 C_*,\qquad R'-\log_3 C_*.
\]
Thus $R'>\log_2C_*+2\Delta$ ensures divisibility by $12^\Delta$.
No assumption that the initial denominator has only primes $2,3$ is
needed: the resulting coordinate is already integral.  The reserves
and the later growth constant $A_0$, used in the source bounds
$E_h=A_0^{h+1}$, are chosen using the new base.

\subsection{Progression counting and the depth-varying extension}

We finish the coefficient check at the finite counting step because
it requires more than the identities for quadratic phases.  In H,
Section~9, eliminating a role $\ell'$ gives the three coefficients
\[
 \alpha_i=i-\ell',\qquad i\in S\setminus\{\ell'\}.
\]
They are nonzero and have absolute value at most four.  The
horizontal frequency tests use sums of role frequencies, while their
anchored tests use the weights $i-\ell'$.  Their kernel cancellations
are precisely \eqref{eq:ll-kernels}.  Accordingly the source's
Fourier expansions, coefficient tails, and annular tests keep their
forms after these replacements.

In H, Lemma~9.1 (\emph{Sparse mass of separated progression tests}),
the shift data consist of a vector $\zeta\in\R^d$ and integer
vectors $n$ with $\|n\|_\infty\le B$.  Replace the shifts
$36n\cdot\zeta$ in its annular relations by $144n\cdot\zeta$,
and the neighborhood shifts $6n\cdot\zeta$ by $12n\cdot\zeta$.
The parameter $N_0$ is the coefficient reserve for the latter
neighborhoods.  The proof uses a fine common grid; choose its step
on a multiple clearing the integers $1,2,3,4$.

For a four-step path, alternately add its four annular relations.
The intermediate label in role $i$ and the two labels in a second
role cancel, leaving $\alpha_i(a-a')$ and an expression in the
four color labels.  Dividing by $\alpha_i$ constructs the source's
difference cover.  The required lattice compatibility is
\begin{equation}
 \frac{144}{\alpha_i}\in12\Z,
 \qquad \frac{12\alpha_j}{\alpha_i}\in\Z.
 \label{eq:ll-counting-divisibility}
\end{equation}
Both assertions hold because $|\alpha_i|$ divides $12$.
The integer vector obtained by adding four original shifts has
norm at most $4B$.  In units of the new neighborhood shift $12$,
division by $\alpha_i$ therefore gives a vector of norm at most
$4B\cdot12/|\alpha_i|\le48B$.
The reserve $N_0\ge64B$ accommodates it.  The scalar error is at most
four times one relation's error after division by a nonzero integer.

For the slab cover, write the first two coordinates as their box
centers plus shifts $12n_i\cdot\zeta$ and scalar interval errors.
Substitute them in the three-role relation and solve for the third
coordinate.  The resulting shift coefficients are
$12\alpha_1/\alpha_3$, $12\alpha_2/\alpha_3$, and
$144/\alpha_3$.  By \eqref{eq:ll-counting-divisibility} they are
integers, and the resulting shift vector has norm at most
$C(d+1)^2B$ as in H.  Choose the separation reserve and
the excluded pure-shift interval constants larger by the fixed
factors arising from these substitutions.  The covering iterations
use the same orders of sumsets, the graph count uses the same four
roles, and the block and scale counts use the same covers.  Their
exponents are consequently unchanged.  In particular the source
factor $L^{29}\delta^{1/2080}$ remains valid, where $L$ bounds the
number of scale occurrences of a fixed pair of labels and $\delta$
is the sparse tuple-density threshold.  In the application, let
$\xi$ be the vector of frequencies obtained from the height
construction and take $\zeta=\xi/144$, so the available heights start at
$12^{\Delta-2}$ instead of $6^{\Delta-2}$.  The preceding height
reserves supply the required shift sizes.  Pair determination in
the application divides only by role differences, and the
short-block windows were checked above.  This verifies the
coefficient-dependent hypotheses of the counting estimate and its
application, not merely its formal conclusion.

We can now apply the remaining operator and summation arguments of
H.  The main projection families, Fourier modes, and counting tests
have the same bounds in powers of logarithms of their budgets.  The
coordinate changes have inserted only fixed rational heights,
fixed derivative losses, and the replacement of one fixed
exponential height base by another.  None changes a budget exponent
as a function of the final $q$.  H's Section~10 chooses those
exponents before taking $q\in(2,3)$ sufficiently close to two;
its absorption of the finite-depth constant therefore applies to
the modified estimates.  This proves \eqref{eq:ll-fixed-local} for
the four role sets together with the intermediate estimates used
in P.

It remains to carry these inputs through P's depth-local proof.
The \emph{Maximal estimate for local catalog rows},
\emph{The depth-varying row estimate}, and
\emph{Smooth mean-zero rows on depth inputs} in P, Section~4,
concern Hilbert-space operators supported on nested dyadic cells.
Here a catalog consists of finitely many vector-valued column
functions, each with a dyadic birth cell; descendants use restrictions
of the same functions.  A local projection is the orthogonal
projection in the cell's $L^2$ space onto constant-coefficient
combinations of a specified set of available columns.
Their hypotheses specify cellwise constant coefficients in a
persistent catalog, an ordinary fixed-input Bessel bound, and
nondecreasing input-read indices.  They contain no role coefficients.
These are P, Lemmas~4.3, 4.4, and 4.6.  The replacement just proved
supplies their same scalar and Bessel inputs.

More concretely, P, Lemma~5.3, Corollary~5.4, and Lemma~5.5,
use
$(y_a,y_b)$ and the demodulations $\pm\gamma/(b-a)$ already checked.
Their cutoff equal to one on $[-1/3,1/3]$ still covers the support
by \eqref{eq:ll-support-t}.  Separating the kernel gives exactly
mean-zero bump rows, with fixed complete carriers.  Clump replacement
and predictor rows have the same Bessel bounds because the
cancellation and height estimates have just been verified.  The
support and observation partitions of the rows, and their
nondecreasing reads of the depth inputs, are unchanged.  In P,
Section~6, use the normalized tags from \eqref{eq:ll-moments}; the
mask, template-transfer, short-block, and counting estimates are
then the versions established above.  Scheduling the column
updates and summing their row bounds introduces no new arithmetic
operation on a role label.

For completeness, the final parameter dependence can be seen in
the retained estimates of H, Section~10.3, and P, Section~7.
Write $k$ for the accuracy logarithm, $\mathfrak p$ for their
growth parameter, $K_0$ for the resulting depth-block length, and
$G_0$ for the strong energy and projection overhead.  The term
requiring $q$ close to two has the form
\[
 \mathfrak p^{-1/2}G_0(K_0+2)^{C_{19}(q-2)},\qquad
 K_0\le C_q(1+k+\mathfrak p)^{C_5},\qquad
 G_0=\exp\!\bigl(C_q(1+\log(2k))^C\bigr).
\]
Here we have renamed the source parameter $P$ to avoid confusion
with the manuscript P.  The degrees $C,C_5,C_{19}$ have absolute
bounds fixed before $q$.  The short-block estimate feeding this
term has the explicit exponent $11(q-2)/(q-1)$ in H, Lemma~7.2;
the coefficient changes above affect its fixed constants, not this
exponent.  Taking $q$ so close to two that
\[
 \max\{C_{19}^{\rm H}C_5^{\rm H},
       C_{19}^{\rm P}C_5^{\rm P}\}(q-2)<1/8
\]
is therefore permitted in both proofs.  Here the accuracy scheme
fixes $0<\theta<1$ before $q$, and the positive constant $\alpha$
is chosen sufficiently small after $q$, as in the source proofs.
Their later choice $\mathfrak p=\lceil\exp(\alpha k^\theta)\rceil$ gives
$G_0=\mathfrak p^{o(1)}$ and makes the displayed term at most
$C_q\mathfrak p^{-1/4}$.  The remaining accuracy exponents and
summations retain their source order of choice.  P's depth-local
proof consequently gives \eqref{eq:ll-depth}.

Translating the role labels back restores $S=\{0\}\cup J$.
P's depth-local theorem permits variation in any role, and the
local constructions allow any role as testing slot.  Thus this
translation does not require the testing role to remain the marked
role $0$ during the local proof.
This completes the proof of Proposition~\ref{prop:ll-depth}.

\subsection{From local forms to orbit oscillation}
\label{subsec:ll-transfer}

An \emph{averaging profile} is a real compactly supported smooth
function $\varphi$ satisfying
\begin{equation}
 \int\varphi=1,\qquad \int t^a\varphi(t)\,\dd t=0\quad(1\le a\le3),
 \qquad
 \|\varphi^{(n)}\|_\infty\le
       (G(n+2))^{G(n+2)}\quad(n\ge0)
 \label{eq:ll-profile}
\end{equation}
for some $G\ge1$.  Profiles may have either sign.  Write
$\varphi_A(t)=A^{-1}\varphi(t/A)$.  These are the profiles of P,
Section~2.

\begin{proposition}[Oscillation for the four triples]
\label{prop:ll-oscillation}
Let $(Y,\nu,R)$ be an invertible probability preserving system,
let $J\subset\{1,2,3,4\}$ have three elements, and let the functions
$g_j:Y\to\C$, $j\in J$, be measurable with $|g_j|\le1$.
For an averaging profile set
\[
 V_k^\varphi(y)=\sum_{n\in\Z}\varphi_{2^k}(n)
                   \prod_{j\in J}g_j(R^{jn}y),\qquad k\ge0.
\]
There is $C_\varphi<\infty$ such that, for all deterministic integers
$0\le k_1<\cdots<k_{W+1}$,
\begin{equation}
 \sum_{b=1}^W\int_Y
   \max_{k_b\le\ell\le k_{b+1}}
      |V_\ell^\varphi-V_{k_b}^\varphi|\,\dd\nu
 \le C_\varphi\sqrt W.
 \label{eq:ll-oscillation}
\end{equation}
The constant is independent of the system, inputs, endpoints, and $W$.
\end{proposition}

\begin{proof}
We give the transfer with the changed coefficients, following P,
Section~3.  Set $b=\varphi_2-\varphi$.  Its moments through order
three vanish, so it has a compactly supported fourth primitive $g$,
with $g^{(4)}=b$.  Choose $B\ge1$ containing its support, a smooth
Gevrey bump $u$ of integral one supported in $[-1/8,1/8]$, and a
dyadic integer $D>32B$.  For $S=\{0\}\cup J$ put
\[
 F(z,\sigma)=u(z)D^{-3}g(D\sigma),\qquad
 \mathcal W(z,\sigma)=
       \prod_{j\in S}(\partial_\sigma-j\partial_z)F(z,\sigma).
\]
The four differential operators commute.  Each is differentiation
along a line on which one role coordinate is fixed, giving all
cancellations in \eqref{eq:ll-kernels}.  On the support,
$|z+j\sigma|\le1/8+4B/D<3/8$.  Thus the rescaled kernel on an
interval of length $r=DA$ has the required interior margin.
The derivative bounds have the form \eqref{eq:ll-kernels}, with
constants depending only on the profile.  Integration in $z$ kills
every term containing a $z$ derivative, leaving
\begin{equation}
 \frac1r\int_\R\mathcal W(z,t/r)\,\dd z
 =A^{-1}b(t/A)=\varphi_{2A}(t)-\varphi_A(t).
 \label{eq:ll-kernel-normalization}
\end{equation}

Let $\Delta_k=\varphi_{2^{k+1}}-\varphi_{2^k}$ and let $f_j$,
$j\in J$, be bounded by one and supported in an interval of length
$L$.  Suppose the testing functions $f_{0,k}$ have the same support
and satisfy the block variation bound $W^{-1/2}$ over a fixed
partition of $k_-\le k\le k_+$ into $W$ blocks.  Apply
Proposition~\ref{prop:ll-depth} to dyadic cells of length $D2^k$,
then average the translated lattice over one period of its largest
scale.  All nonzero cells lie below roots of total length at most
$L+2D2^{k_+}$.  Scale order is the reverse of depth order, which
preserves block variation.  The normalized translation average
of a scale's kernels is \eqref{eq:ll-kernel-normalization}.
We obtain
\begin{equation}
 \left|\sum_{k=k_-}^{k_+}\int_{\R^2}
     f_{0,k}(x)\prod_{j\in J}f_j(x+jt)\Delta_k(t)
       \,\dd x\,\dd t\right|
 \le C_\varphi(L+D2^{k_+}).
 \label{eq:ll-continuous-test}
\end{equation}

For clarity, integer transfer does not require consecutive roles.
Plant each integer sequence on intervals of radius $e<1/200$ about
the integers.  If $x+jt$ lies in such intervals in all four roles,
use an adjacent pair to reconstruct integers $m,n$.  Every other
integer center must equal $m+jn$, since its discrepancy has
magnitude less than $8e$.  Write $x=m+\xi$, $t=n+\eta$.
The offset domain is
\[
 \Omega_{S,e}=\{(\xi,\eta):|\xi+j\eta|<e\text{ for all }j\in S\}.
\]
If $s_-=\min S$ and $s_+=\max S$, the intermediate inequalities
follow from the two extreme ones by convexity.  Hence
\begin{equation}
 |\Omega_{S,e}|=\frac{4e^2}{s_+-s_-}.
 \label{eq:ll-polygon}
\end{equation}
Replacing $\Delta_k(n+\eta)$ by $\Delta_k(n)$ has, after summing
in $n$, error $O_\varphi(2^{-k})$ per integer $m$: its derivative
is $O_\varphi(2^{-2k})$, it has $O_\varphi(2^k)$ relevant integer
arguments, and $|\eta|\le2e$.  Summing over $k\ge0$ costs
$O_\varphi(L)$.  Dividing by the fixed area
\eqref{eq:ll-polygon} therefore gives the discrete counterpart of
\eqref{eq:ll-continuous-test} with bound
$C_\varphi(L+D2^{k_+})$.

Fix a finite interval $\mathcal E$ of integer testing points
inside the common support interval of the integer sequences.
For each $m\in\mathcal E$ and block $b$, choose a maximizing
endpoint $\ell_b(m)$ and a unimodular scalar $\alpha_b(m)$ making
the selected difference nonnegative.  Set, on that scale block,
\[
f_{0,k}(m)=\frac{\alpha_b(m)}{4\sqrt W}
       \mathbf1_{\{k_b\le k<\ell_b(m)\}}.
\]
Set $f_{0,k}(m)=0$ for $m\notin\mathcal E$.
Its supremum plus variation is at most $1/(2\sqrt W)$.
The discrete test telescopes into $1/(4\sqrt W)$ times the sum
of the block maxima.  Thus its oscillation bound is
$C_\varphi\sqrt W(L+D2^{k_{W+1}})$.

We transfer the sequence inequality along finite orbit segments,
following Calder\'on's transference method \cite{Calderon}.
Apply the estimate to the finite orbit sequences
$g_j(R^r y)$ for $1-Q\le r\le M+Q$, set to zero elsewhere, where
$Q=\lceil4B_\varphi2^{k_{W+1}}\rceil$ and
$\supp\varphi\subset[-B_\varphi,B_\varphi]$.
Take $\mathcal E=\{1,\ldots,M\}$.  At these testing points all
averages are the untruncated
orbit averages evaluated at $R^m y$.  Integrate in $y$, use
invariance, divide by $M$, and let $M\to\infty$.  The result is
\eqref{eq:ll-oscillation}.  This step uses only invertibility,
measurability, and preservation of the probability measure.
\end{proof}

\subsection{Convergence on arbitrary systems}
\label{subsec:ll-convergence}

We now prove Lemma~\ref{lem:lower-length} from
Proposition~\ref{prop:ll-oscillation}.  The use of quantitative
oscillation bounds to prove pointwise convergence belongs to the
classical approach developed in \cite{JonesKaufmanRosenblattWierdl}.
The criterion needed here has the following short proof.
A uniformly bounded
sequence of measurable functions satisfying
\eqref{eq:ll-oscillation} is Cauchy almost everywhere.  Indeed, if
its tail oscillation exceeds $4\epsilon$ on a set $E$ of positive
measure, then for every fixed initial endpoint $k$ the supremum
of $|V_\ell-V_k|$ over $\ell\ge k$ exceeds $2\epsilon$ on $E$.
By monotone convergence one can choose a deterministic later endpoint
so that the corresponding integral maximum is at least
$\epsilon\nu(E)$.  Recursively choosing endpoints gives a lower
bound $W\epsilon\nu(E)$, contradicting the upper bound
$C_\varphi\sqrt W$ for large $W$.  The finite sums defining
$V_k^\varphi$ are uniformly bounded because the profile is bounded
and compactly supported.  Thus every fixed profile gives almost
everywhere convergence.

P, Lemma~2.8 (\emph{Approximation with three vanishing moments}),
states that for every $c\in[1,2]$ and every $\epsilon>0$ there is
an averaging profile with
\begin{equation}
 \|\varphi-h_c\|_1<\epsilon,
 \qquad h_c(t)=c^{-1}\mathbf1_{(0,c]}(t).
 \label{eq:ll-profile-approximation}
\end{equation}
The use of signed profiles is essential to this assertion.  One
first smooths $h_c$, preserving integral one.  The other three
moments can be canceled, without changing the integral, by four
broad translated bumps at distance and width comparable to a large
parameter $A$.  Their moment matrix for orders $0,1,2,3$, after
division of row $j$ by $A^j$, is a fixed invertible Vandermonde-type
matrix.  The prescribed zeroth correction is zero and the remaining
normalized corrections are $O(A^{-1})$.  The correcting coefficients
are therefore $O(A^{-1})$, and their total $L^1$ cost tends to zero.  The bumps
retain the derivative bounds in \eqref{eq:ll-profile} with a
profile-dependent constant.

Fix a triple $J$ and write
$a_n(y)=\prod_{j\in J}g_j(R^{jn}y)$, so $|a_n(y)|\le1$.
For $N_k=\lfloor c2^k\rfloor$, the Riemann-sum comparison gives,
uniformly in $y$,
\begin{equation}
 \limsup_{k\to\infty}
 \left|\frac1{N_k}\sum_{n=1}^{N_k}a_n(y)-V_k^\varphi(y)\right|
 \le\|h_c-\varphi\|_1.
 \label{eq:ll-profile-error}
\end{equation}
To see this, first replace $N_k^{-1}$ by $(c2^k)^{-1}$, at cost
$O(2^{-k})$, then bound the difference by
$2^{-k}\sum_{n\in\Z}|h_c(n2^{-k})-\varphi(n2^{-k})|$.
This converges to the displayed integral because the integrand is
compactly supported and piecewise continuous.

Choose countably many profiles with errors tending to zero for each
rational $c\in[1,2]$.  On their common conull set, the existence of
all the smooth-profile limits and \eqref{eq:ll-profile-error} show
that the ordinary averages along $N_k$ are Cauchy.  The
Host--Kra norm convergence theorem \cite[Theorem~1.1]{HostKra},
stated for arbitrary invertible probability preserving systems,
gives an $L^2$ limit for the averages at all integer lengths: apply
it to the four consecutive powers, inserting the constant function
one at the omitted role.  The bounded pointwise limit along each
$N_k$ must equal this $L^2$ limit almost everywhere, by bounded
convergence and uniqueness of norm limits.  Thus all rational-$c$
subsequences have the same limit on one conull set.  No identification
of that limit and no mixing assumption is needed.

Finally, for $M\le N$ and $|a_n|\le1$,
\[
 \left|\frac1N\sum_{n=1}^Na_n-\frac1M\sum_{n=1}^Ma_n\right|
 \le 2\frac{N-M}{N}.
\]
For any rational mesh of $[1,2]$, every sufficiently large integer
$N$ lies within its relative mesh size, up to $O(N^{-1})$, of some
$\lfloor c2^k\rfloor$ from that finite mesh.  Its finitely many
subsequences have the common pointwise limit.  The preceding
inequality, followed by meshes of size tending to zero, proves
convergence along all positive integers.  Homogeneity removes the
normalization $|g_j|\le1$.  A proper subset with fewer than three
roles is contained in one of the triples already treated; assign
the constant function one to its additional roles.  This proves
Lemma~\ref{lem:lower-length} in its full stated scope.
